\documentclass[11pt]{article}
\usepackage{manual_docs_style}
\externaldocument{build/environment_profiles}[environment_profiles.pdf]
\externaldocument{build/run_graph_schema}[run_graph_schema.pdf]
\externaldocument{build/simulation_stage}[simulation_stage.pdf]

\begin{document}

\docTitle{Stage: Compute Metrics}
\phantomsection\label{docstart:compute_metrics_stage}

This stage computes detectability and eligibility metrics from materialized simulation runs.
The result is used by dataset selection and question generation to reject impossible or ambiguous samples.

The main driver is:

\begin{DocCode}
env/bin/python workflows/metrics/compute_metrics.py [--model MODEL] --policy POLICY_ID [--jobs N]
\end{DocCode}

\section*{Inputs}

For each model, the stage reads:

\begin{itemize}
\item \code{plans/<policy_id>/run_nodes.jsonl}: canonical run recipes and \code{run_id} values.
\item \code{plans/<policy_id>/run_edges.jsonl}: baseline/intervention/time-zero relationships.
\item \code{runs/model_record.json}: runtime status and recipe hashes.
\item \code{runs/<run_id>/data.parquet}: materialized trajectories.
\item \code{experiment_config.json}: observed signals, signal types, and detectability thresholds.
\item \code{noise_adder.py}: model-specific noise profiles and SNR analysis.
\item \code{detectability_specific_environment.py}: optional environment-specific detectability veto hook.
\end{itemize}

During migration, the stage may also read legacy \code{model_run_specs.json}, but new code should use the resolved run graph.

\section*{Run Relationships}

For each intervention node, the stage uses \code{run_edges.jsonl} to find:

\begin{itemize}
\item the matched no-intervention baseline run;
\item the matched time-zero/static-change probe run;
\item the changed root-cause parameter;
\item the intervention time and post-intervention value.
\end{itemize}

These relationships are not inferred from filename or identifier structure.

\section*{Detectability Computation}

For each successful intervention run, compute detectability against:

\begin{itemize}
\item its matched baseline run;
\item its matched time-zero/static-change run.
\end{itemize}

Both comparisons load the dataframes, assert aligned timestamps, apply model-specific noise when requested, preprocess channels according to the environment and paper appendix, compute SRD threshold crossings, compute RMS distances, and compute effect-to-noise SNR values.

The generic output includes:

\begin{itemize}
\item \code{mean_rms_distance_clean_dirty}
\item \code{mean_rms_distance_clean_baseline}
\item \code{mean_SNR}
\item \code{first_diff}
\item \code{detectable}: \code{"yes"}, \code{"no"}, or \code{"error"}
\end{itemize}

For the comparison against the no-intervention baseline, the stage may additionally call:

\begin{DocCode}
detectability_specific_environment.py::is_detectable(...)
\end{DocCode}

If this hook returns \code{"no"}, the intervention is not detectable even if the generic RMS/SRD checks pass.

\section*{Eligibility Rules}

An intervention is eligible only if:

\begin{itemize}
\item it has status \code{success} in \code{model_record.json};
\item its comparison against the matched baseline has \code{detectable == "yes"};
\item its comparison against the matched time-zero/static-change run has \code{detectable == "yes"}.
\end{itemize}

A family is eligible when it has eligible children for at least the configured minimum number of distinct root-cause parameters.
The default minimum is three distinct parameters unless a policy or selection stage specifies a stricter rule.

\section*{Output}

The preferred output location is:

\begin{DocCode}
models/simulink/<MODEL>/metrics/<policy_id>/eligibility_metrics.jsonl
\end{DocCode}

A legacy aggregate file may also be written during migration:

\begin{DocCode}
models/simulink/<MODEL>/runs/eligibility_metrics.json
\end{DocCode}

\section*{Preferred JSON Lines Shape}

Each line of \code{eligibility_metrics.jsonl} records either one intervention or one family summary.

Example intervention record:

\begin{DocCode}
{
  "record_type": "intervention",
  "policy_id": "benchmark_v1",
  "family_id": "fam_de4d7c0a5335e9005510cc2442d0971e",
  "run_id": "run_7fd9daa0055fb3f412ef5e0ab8d4c5ba",
  "baseline_run_id": "run_baseline_...",
  "time0_run_id": "run_time0_...",
  "changed_parameter": "gravity",
  "eligible": true,
  "detectability": {
    "vs_baseline": {"detectable": "yes"},
    "vs_time0_baseline": {"detectable": "yes"}
  }
}
\end{DocCode}

Example family record:

\begin{DocCode}
{
  "record_type": "family",
  "policy_id": "benchmark_v1",
  "family_id": "fam_de4d7c0a5335e9005510cc2442d0971e",
  "baseline_run_id": "run_baseline_...",
  "family_eligible": true,
  "eligible_parameters": ["gravity", "mass", "drag_coeff"]
}
\end{DocCode}

\section*{Validation Rules}

\begin{itemize}
\item Every metric record must refer to a \code{run_id} present in \code{run_nodes.jsonl}.
\item Every intervention metric must have a matched baseline and time-zero/static-change run from \code{run_edges.jsonl}.
\item Only successful materialized runs may be marked eligible.
\item The metric stage must not change run recipes or recipe hashes.
\item Re-running this stage may overwrite previous metrics for the same \code{policy_id}.
\end{itemize}

\end{document}
